\subsection{Phases of a Compiler}

The phases are:
\begin{itemize}
    \item scanner
    \item parser 
    \item semantic actions
    \item optimizer
    \item code generator
\end{itemize}

The compiler starts by seeing only text,
\texttt{'i' 'f' ' ' '(' ...}. The \emph{scanner}
phase converts raw text to a series of tokens, 
one for each "word" in the program, 
\texttt{if ( ID(a) OP(<) LIT(4) ) ...}.

The \emph{parser} then takes this string of 
tokens and coverts it to a \emph{parse tree} or 
\emph{abstract syntax tree} which captures the 
structure of the code, "this is an if-then, 
in the first case follow this code path...". 

Syntax tells what symbols, in what order, 
are legal in a language. Semantics 
is the meaning. The phrase "gorilla yip 
huffly huffly" is legal syntax in 
English, but has no semantic meaning. 
Semantic actions are taken by compiler 
based on the semantics of program statements,
including building a \emph{symbol table} and
generating \emph{intermediate representations}.

Symbol tables associate each identifier or 
symbol (variables, constants, 
procedures, functions, etc.) in a program's 
code with the necessary information for its 
functioning. This may include type and scope of a variable,
return type for a function, etc. 

Intermediate representation is a
low level representation of the program that 
is not machine specific (and therefore does 
not reference registers), such as 
three-address code. 

The \emph{optimizer} transforms code to make 
it more efficient. There are many optimizations 
to make, especially for code produced by students. 

The \emph{code generator} generates assembly from 
the intermediate represention. 